\chapter{\textbf{Materiais e Métodos}}

A seguir s...

\section{Materiais e Reagentes}

\subsection{Materiais}

 \begin{itemize}
   \item Reator de bancada com volume de 5L;
   
   \item Refratômetro portátil.
   
 \end{itemize}

\subsection{Reagentes}

 \begin{itemize}
   \item Cloreto benzalcônio 50 \%;
   \item Meio de cultura Luria Bertani (LB).
   
 \end{itemize}

\section{Métodos}

A metodologia ...

\subsection{Levantamento de Mercado}
\label{sub:levantamentodemercado}

Realizou-se u..

%EXEMPLO DE EQUAÇÃO

Para o cálculo dos volumes necessários de cloreto benzalcônio 50 \% (v/v) (CB50) e de álcool etílico 92,8 \% (v/v) para obtenção das concentrações finais desejadas foi utilizada a Equação \ref{equa:diluicao}.


\begin{gather}
	\label{equa:diluicao} 
    C_{inicial} \ . \ V_{inicial} \ = \ C_{final} \ . \ V_{final} 
\end{gather}

\noindent em que,\\
 $C_{inicial}$: Concentração inicial, antes da diluição (\%); \\
 $V_{inicial}$: Volume inicial, antes da diluição (mL); \\
 $C_{final}$: Concentração final (\%); \\
 $V_{final}$: Volume final (mL). \\
 





